\documentclass[11pt]{article}
\usepackage[T1]{fontenc}
\usepackage{lmodern}
\usepackage[margin=1in]{geometry}
\usepackage{amsmath,amssymb,amsthm,mathtools}
\usepackage{microtype,booktabs,longtable,array,enumitem}
\usepackage{xcolor}
\usepackage[colorlinks=true,linkcolor=blue!45!black,citecolor=blue!45!black,urlcolor=blue!45!black]{hyperref}
\usepackage{fancyhdr}
\pagestyle{fancy}\fancyhf{}\fancyhead[L]{\small Complementary depth in polylogarithms}\fancyhead[R]{\small ProveIt research contribution}\fancyfoot[C]{\thepage}
\setlength{\headheight}{14pt}
\setlength{\emergencystretch}{2em}
\numberwithin{equation}{section}
\newtheorem{theorem}{Theorem}[section]
\newtheorem{proposition}[theorem]{Proposition}
\newtheorem{corollary}[theorem]{Corollary}
\newtheorem{lemma}[theorem]{Lemma}
\theoremstyle{definition}\newtheorem{definition}[theorem]{Definition}
\theoremstyle{remark}\newtheorem{remark}[theorem]{Remark}
\DeclareMathOperator{\Li}{Li}
\DeclareMathOperator{\Rea}{Re}
\DeclareMathOperator{\Ima}{Im}
\DeclareMathOperator{\dep}{dep}
\DeclareMathOperator{\wt}{wt}
\newcommand{\Q}{\mathbb Q}\newcommand{\C}{\mathbb C}
\newcommand{\emp}{\varnothing}\newcommand{\sh}{\mathbin{\amalg}}
\newcommand{\ones}[1]{\{1\}^{#1}}
\newcommand{\code}[1]{\texttt{\detokenize{#1}}}
\title{Complementary depth, classical ladders,\texorpdfstring{\\}{ }and exact shuffle ranks\large\medskip\\Gaussian and Eisenstein polylogarithms:\\proofs of three manuscript candidates and a mixed-point correction}
\author{Research draft prepared for Vladimir Reshetnikov\\\small ProveIt / Analysis / Polylogarithms}
\date{9 October 2026}
\begin{document}
\maketitle
\begin{abstract}
We develop a branch-explicit reduction for one-variable multiple
polylogarithms at the transformed argument $q=(1-z)^{-1}$. A word of weight
$w$ and depth $d$ is expressed as a finite rational combination of
$\log^j q\,\Li_{\boldsymbol a}(q)\zeta(\boldsymbol b)$ with
$\dep(\boldsymbol a)+\dep(\boldsymbol b)\le w-d$. The logarithm is a
coefficient in this depth statement; the argument family changes.
The one-zero specialization yields classical formulas for every
$\Li_{\ones a,2,\ones b}(z)$, and proves all three weight-four Gaussian
triple identities printed as numerical candidates in the inspected
manuscript. A particularly short terminal-$2$ ladder and an all-weight
double-polylogarithm formula for $\Li_{3,\ones b}(z)$ give further explicit
identities, including the Eisenstein specialization.

Separately, a multigraded Lyndon-word calculation determines the exact
formal shuffle-product rank at every weight and depth. It explains the
reported rank seven on ten weight-six triples without implying numerical
independence. We correct an incorrect mixed-point antisymmetry statement
by an exact nonzero Gaussian diagonal example. We also record a rational
dual witness showing that the manuscript's $S_4$ target is outside a
specified finite depth-two double-shuffle row space; $S_4$ itself remains
unproved here. Exact word certificates, 153 independent quadrature
checks, and rational interval enclosures accompany the proofs.
\end{abstract}
\vspace{3pt}
\noindent\textbf{Status.} Unrefereed research contribution with self-contained
analytic arguments and reproducible finite checks. Not proof-assistant
verified. The general harmonic-polylogarithm and Lyndon-algebra machinery
is established mathematics; no global priority claim is made for it.
\clearpage
\tableofcontents
\clearpage

\section{Scope, provenance, and the mathematical contribution}

The input is Chapter 4, \emph{Gaussian and Eisenstein depth}, of the
ProveIt polylogarithm manuscript, together with its discovery chapter
\cite{ProveIt}. We inspected the repository at commit
\begin{center}\small\code{29d9344771b41df39ff9b7890c76b1f152572dde}.\end{center}
The Chapter 4 blob is
\code{6171912ff7724784cabb789a066c61223ab298bc}.
These identifiers distinguish this contribution from subsequent changes
to a rapidly evolving repository.

The manuscript lists three formulas for
$\Ima\Li_{2,1,1}(i,1,1)$,
$\Ima\Li_{1,2,1}(i,1,1)$, and
$\Ima\Li_{1,1,2}(i,1,1)$ among numerically supported Gaussian reductions.
It also reports a shuffle rank of seven on ten weight-six triple
coordinates, and asks how a formal quotient relates to actual depth.
These are the principal targets of this article. The existing
weight-six double parity formulas and inverse-color gap reductions are
not the main subject of the present continuation.

The results have three logically different forms. The functional
identities are equalities of analytic functions on an explicitly stated
domain and hence equalities of their special values. The shuffle-rank
result is an exact dimension theorem for a specified \emph{formal}
algebra. The numerical calculations are either diagnostics or
proof-backed enclosures; neither kind implies a new independence theorem.
Keeping those distinctions visible is essential here.

Harmonic polylogarithms, their change-of-variable formulas, and the
Nielsen subfamilies have a substantial classical literature
\cite{RemiddiVermaseren,Kolbig}. The structural shuffle-algebra theorem
used below is due to Radford, with a useful triangular formulation in
Melan\c{c}on--Reutenauer \cite{Radford,MelanconReutenauer}.
Our contribution is the explicit branch-controlled organization needed
for these manuscript questions: closed formulas with all endpoint
constants retained, exact specializations of the printed candidates,
a finite certificate catalogue, and corrections that prevent the formal
calculations from being misread as arithmetic lower bounds.

\subsection{A map of the principal results}
The complementary-depth theorem is Theorem~\ref{thm:complement}.
The complete one-zero reduction and terminal-$2$ ladder are
Theorems~\ref{thm:onezero} and~\ref{thm:terminal}.
The manuscript's three Gaussian candidates become
Theorem~\ref{thm:gaussian}. The two-zero height-one family is
Theorem~\ref{thm:twozero}. The all-weight formal rank is
Theorem~\ref{thm:rank}. Proposition~\ref{prop:mixedcorrection} gives the
nonzero diagonal counterexample. None of these statements asserts
numerical independence of the weight-six triple candidates.

\section{Conventions and analytic domains}

For positive indices we abbreviate
\begin{equation}\label{eq:mpl}
\Li_{s_1,\ldots,s_d}(z)
 :=\Li_{s_1,\ldots,s_d}(z,1,\ldots,1)
 =\sum_{n_1>\cdots>n_d\ge1}
 \frac{z^{n_1}}{n_1^{s_1}\cdots n_d^{s_d}},\qquad |z|<1.
\end{equation}
The indices decrease in the summation order, as in the manuscript.
Weight is $w=s_1+\cdots+s_d$ and series depth is $d$.
An empty index denotes the multiplicative identity.
We use convergent multiple zeta values
$\zeta(\boldsymbol s)=\Li_{\boldsymbol s}(1)$ when $s_1\ge2$.

Let
\[
 \omega_0(t)=\frac{dt}{t},\qquad
 \omega_1(t)=\frac{dt}{1-t}.
\]
Words are written \emph{outer letter first}. Thus
\begin{equation}\label{eq:word}
H_{0^{s_1-1}1\cdots0^{s_d-1}1}(z)=\Li_{s_1,\ldots,s_d}(z).
\end{equation}
In particular $H_{01}=\Li_2$, not its negative or a reversed word.
The recursion is
$H_{av}(z)=\int_0^z\omega_a(t)H_v(t)$ whenever this is an ordinary
convergent integral. Pure-zero words are defined by
$H_{0^r}(z)=\log^r z/r!$, and trailing-zero words are obtained by shuffle
regularization. Every word ending in $1$ is convergent at its initial
endpoint and has the ordinary interpretation in \eqref{eq:word}.

Throughout the functional identities we use
\begin{equation}\label{eq:domain}
 D=\C\setminus[0,\infty),\qquad
 q=\frac1{1-z},\qquad L=\log q=-\log(1-z),\qquad z\in D,
\end{equation}
with principal logarithms and principal polylogarithms. The image of $D$
avoids both $(-\infty,0]$ and $[1,\infty)$ in the $q$-plane, so the
hyperlogarithms with possible trailing zeros have unambiguous branches.
The origin is a limiting point, not an included point of $D$.

We first prove formulas on $z<0$, where $0<q<1$ and all logarithms in
$q$ are real. Analytic continuation on the connected slit domain $D$
then proves the stated complex identities. This method avoids an
unstated choice of a value on the polylogarithm cut. It includes both
$z=i$ and $z=\rho=e^{2\pi i/3}$:
\begin{align}
 z=i:&\quad q=\frac{1+i}{2},&L&=-\frac12\log2+\frac{i\pi}{4},\label{eq:gausspoint}\\
 z=\rho:&\quad q=\frac12+\frac{i\sqrt3}{6},&L&=-\frac12\log3+\frac{i\pi}{6}.\label{eq:eispoint}
\end{align}
Their transformed moduli are $2^{-1/2}$ and $3^{-1/2}$ respectively.

\begin{definition}[Depth with logarithmic coefficients]
A term $L^j\Li_{\boldsymbol a}(q)\zeta(\boldsymbol b)$ has
\emph{reduced depth} $\dep(\boldsymbol a)+\dep(\boldsymbol b)$.
The factor $L^j$ is a coefficient, not an additional nested sum.
\end{definition}
This is an explicit convention for the following reduction. It is not
an identification with every motivic depth filtration, and does not
turn arbitrary products of deep factors into depth-one objects.

\section{An explicit complementary-depth theorem}

\subsection{Endpoint transport with convergent constants}
For a word $v$ ending in $0$, write $J_v(q)$ for the iterated integral
from $1$ to $q$ using the same forms and the same outer-first convention.
These integrals converge at $1$: the innermost form is $\omega_0$, which
is regular there; preceding $\omega_1$ integrations do not destroy
integrability. Set $J_\emp=1$.

\begin{lemma}[Connection formula]\label{lem:connection}
For every nonempty word $v$ ending in $0$ and $0<q<1$,
\begin{equation}\label{eq:connection}
 J_v(q)=\sum_{v=uv'}(-1)^{|v'|}H_u(q)H_{\operatorname{rev}(v')}(1).
\end{equation}
The sum includes empty prefixes and suffixes. Every nonempty endpoint
word $\operatorname{rev}(v')$ starts in $0$, so its value at $1$ is
convergent after the usual trailing-zero normalization at $0$.
\end{lemma}
\begin{proof}
We give the endpoint justification rather than silently substituting into
a divergent path identity. Fix a maximum word length $N$ and work in the
truncated noncommutative algebra on $e_0,e_1$. Let
$F(q)=\sum_{|u|\le N}H_u(q)u$. It satisfies
\[
 F'(q)=\left(\frac{e_0}{q}+\frac{e_1}{1-q}\right)F(q).
\]
Put $\ell=-\log(1-q)$. The function $e^{-\ell e_1}F(q)$ has a finite
coefficientwise limit $Z$ as $q\to1^-$: its differential equation has
coefficient $e^{-\ell e_1}e_0e^{\ell e_1}/q$, a finite polynomial in
$\ell$ at each truncated weight, and this is integrable at $1$.
Consequently
$F(q)=e^{\ell e_1}Z+O((1-q)\log^M(1-q))$
coefficientwise, for a sufficiently large finite $M$.

Shuffle identities say that $F$ is group-like; the same holds for $Z$.
The inverse coefficient in a group-like word series is
$(Z^{-1})_{v'}=(-1)^{|v'|}Z_{\operatorname{rev}(v')}$, by the shuffle
antipode, or equivalently by reversal of an iterated path.
The coefficients of $R(q)=F(q)Z^{-1}$ solve the same differential
recursion and tend to zero at $1$ for every nonempty word ending in $0$.
Induction on word length therefore identifies them with the convergent
integrals $J_v(q)$.

If $v'$ is a nonempty suffix of $v$, its reverse begins in $0$.
The corresponding coefficient of $Z$ is simply
$H_{\operatorname{rev}(v')}(1)$: multiplication by
$e^{-\ell e_1}$ only changes words beginning with $1$.
Taking the coefficient of $v$ in $F(q)Z^{-1}$ proves
\eqref{eq:connection}. Since $N$ was arbitrary, this proves every finite
word identity.
\end{proof}

\subsection{The word formula and the depth bound}
For a word $w$ ending in $1$, let $K$ be its number of zeros. Replace
each original $1$ by $0$, and replace each original $0$ independently
by either $0$ or $1$. Denote the resulting $2^K$ words by $\tau_A(w)$,
where $A$ records the zero positions replaced by $1$.
Every $\tau_A(w)$ ends in $0$.

\begin{theorem}[Complementary depth]\label{thm:complement}
For $z\in D$ and every nonempty word $w$ ending in $1$,
\begin{equation}\label{eq:master}
 H_w(z)=(-1)^K\sum_A\ \sum_{\tau_A(w)=uv}
 (-1)^{|v|}H_u(q)H_{\operatorname{rev}(v)}(1).
\end{equation}
After trailing-zero normalization, this is a finite rational linear
combination of terms
\begin{equation}\label{eq:termshape}
 L^j\Li_{\boldsymbol a}(q)\zeta(\boldsymbol b),\qquad
 j+\wt(\boldsymbol a)+\wt(\boldsymbol b)=|w|,
 \quad \dep(\boldsymbol a)+\dep(\boldsymbol b)\le K.
\end{equation}
Nonempty zeta indices are convergent. Thus a one-variable polylogarithm
of weight $w$ and depth $d$ has the stated reduction with $K=w-d$.
\end{theorem}
\begin{proof}
On $z<0$, make the substitution $u=(1-t)^{-1}$ in every integration.
Its pullbacks are
\begin{equation}\label{eq:pullbacks}
 \frac{dt}{t}=-\frac{du}{u}-\frac{du}{1-u},\qquad
 \frac{dt}{1-t}=\frac{du}{u}.
\end{equation}
The initial endpoint $t=0$ becomes $u=1$. Multilinearity gives
$H_w(z)=(-1)^K\sum_A J_{\tau_A(w)}(q)$.
Lemma~\ref{lem:connection} proves \eqref{eq:master}.

It remains to check the claimed normal form. The identity
$H_vH_0=H_{v\sh0}$ removes a trailing zero. If a word ends in exactly
$r$ zeros, its coefficient in the shuffle of its prefix with $0$ is
$r$; all other words in that shuffle have fewer trailing zeros.
Induction therefore expresses each $H_u(q)$ as a rational combination
of $L^jH_{u'}(q)$ with $u'$ empty or ending in $1$.
The operation preserves the number of $1$ letters.

For an endpoint word beginning in $0$, the same normalization preserves
its initial $0$ in every nonzero term with no logarithmic factor.
Terms containing a positive power of $\log1$ vanish. The surviving
nonempty terms begin in $0$ and end in $1$, hence are convergent MZVs.
In each split of $\tau_A(w)$ the total number of $1$ letters in the
two factors is $|A|\le K$. This proves the depth bound and the weight
identity. Finally continue analytically from $z<0$ to $D$.
\end{proof}

\begin{corollary}[An argument-enlarged depth bound]\label{cor:halfb}
For $w\ge2$, every one-variable MPL of weight $w$ can be represented
using nonlogarithmic factors of depth at most $\lfloor w/2\rfloor$,
allowing both the original argument $z$ and $q=(1-z)^{-1}$, MZVs, and
powers of $L$ as coefficients. More precisely, use the original
expression when $d\le w/2$, and Theorem~\ref{thm:complement} otherwise.
The all-one case is $L^w/w!$.
\end{corollary}

This corollary is deliberately not phrased as a theorem about arbitrary
multivariable MPLs, or as a reduction inside an unchanged root-of-unity
basket. At $z=i$ the half-point $q$ is not a root of unity. The theorem
explains an argument-enlarged upper bound; it does not prove the absence
or presence of new cyclotomic periods in a smaller comparison algebra.

\section{The complete one-zero classical ladder}

The one-zero case admits a simpler derivation that is independent of
the endpoint-transport implementation.
For $r\ge1$ and $a,b\ge0$, put
\[
 F_r(z)=\Li_{2,\ones{r-1}}(z),\qquad
 T_{a,b}(z)=\Li_{\ones a,2,\ones b}(z).
\]
The word of $T_{a,b}$ is $1^a01^{b+1}$.

\begin{theorem}[One-zero reduction]\label{thm:onezero}
On $D$,
\begin{align}
 F_r(z)
 &=\sum_{k=1}^{r+1}\frac{(-1)^kL^{r+1-k}}{(r+1-k)!}\Li_k(q)
   +(-1)^r\zeta(r+1)-\frac{L^{r+1}}{(r+1)!},\label{eq:Fr}\\
 T_{a,b}(z)
 &=\sum_{j=0}^{a}(-1)^j\binom{b+j+1}{j}
    \frac{L^{a-j}}{(a-j)!}F_{b+j+1}(z).\label{eq:Tab}
\end{align}
In particular every such MPL, however large its series depth, belongs
to the algebra generated by $L$, ordinary zeta values, and classical
polylogarithms $\Li_k(q)$ of weight at most $a+b+2$.
\end{theorem}
\begin{proof}
Repeated integration of the identical $\omega_1$ letters gives
\begin{equation}\label{eq:Frintegral}
 F_r(z)=\frac1{r!}\int_0^z\frac{[-\log(1-t)]^r}{t}\,dt.
\end{equation}
For $z<0$ substitute $u=(1-t)^{-1}$ to obtain
\[
 F_r(z)=\frac1{r!}\int_1^q\frac{\log^r u}{u(u-1)}\,du.
\]
A primitive of $\log^r u/(1-u)$ is
\[
 A_r(u)=\sum_{k=1}^{r+1}
 (-1)^{k-1}\frac{r!}{(r+1-k)!}\log^{r+1-k}u\,\Li_k(u).
\]
Differentiation telescopes. At $u=1$ only the last term survives, with
$A_r(1)=(-1)^rr!\zeta(r+1)$; the apparent
$\log^r u\,\Li_1(u)$ singularity tends to zero because $r\ge1$.
Using $1/[u(u-1)]=-1/u-1/(1-u)$ proves \eqref{eq:Fr}.

The $a$ outer $\omega_1$ integrations can be performed first:
\begin{equation}\label{eq:singleintegral}
 T_{a,b}(z)=\frac1{a!(b+1)!}\int_0^z
  \bigl(L+\log(1-t)\bigr)^a[-\log(1-t)]^{b+1}\frac{dt}{t}.
\end{equation}
Expand the first power and apply \eqref{eq:Frintegral}.
The resulting factorial coefficient is exactly the binomial coefficient
in \eqref{eq:Tab}. The integrals are convergent at $0$, and continuation
on $D$ completes both proofs.
\end{proof}

\subsection{A particularly short \texorpdfstring{terminal-$2$}{terminal-2} identity}
When the $2$ is the last index, substantial cancellation leaves only two
classical polylogarithms at $q$.

\begin{theorem}[Terminal-$2$ ladder]\label{thm:terminal}
For every integer $a\ge0$,
\begin{align}
\Li_{\ones a,2}(z)
={}&-\frac{L^{a+2}}{(a+2)!}
 -\sum_{m=2}^{a+1}\frac{(m-1)L^{a+2-m}}{(a+2-m)!}\zeta(m)\notag\\
 &-L\Li_{a+1}(q)+(a+1)\Li_{a+2}(q)
 -(a+1)\zeta(a+2).\label{eq:terminal}
\end{align}
An empty sum is zero.
\end{theorem}
\begin{proof}
For $a=0$ this is \eqref{eq:Fr} with $r=1$.
Differentiate the right side with respect to $L$, using $q=e^L$.
For $a\ge1$ the result is the right side for $a-1$; in particular the
last term of the finite zeta sum supplies the required constant
$-a\zeta(a+1)$. The same recursion holds on the left because an outer
$\omega_1$ integration is integration in $L$.
Both sides tend to zero as $q\to1^-$, or $z\to0^-$.
The term $L\Li_1(q)$ tends to zero in the base case.
Induction and continuation prove the formula.
\end{proof}

For example, two explicit higher-depth identities are
\begin{align}
\Li_{1,1,1,2}(z)
={}&4\Li_5(q)-L\Li_4(q)-4\zeta(5)-3L\zeta(4)
-L^2\zeta(3)\notag\\[-2pt]
&-\frac{L^3}{6}\zeta(2)-\frac{L^5}{120},\label{eq:new5}\\
\Li_{1,1,1,1,2}(z)
={}&5\Li_6(q)-L\Li_5(q)-5\zeta(6)-4L\zeta(5)
-\frac32L^2\zeta(4)\notag\\[-2pt]
&-\frac13L^3\zeta(3)-\frac1{24}L^4\zeta(2)-\frac{L^6}{720}.
\label{eq:new6}
\end{align}
These are proved functional identities, not PSLQ conjectures. Either
Gaussian or Eisenstein substitution gives an explicit family of new
entries for the manuscript's tables.

\section{Proof of the three Gaussian triple candidates}

Write $\ell=\log2$, $G=\beta(2)$, and
$\lambda_n=\Ima\Li_n((1+i)/2)$ as in the manuscript.
The low-weight classical evaluations needed below are
\begin{align}
\Li_1(q)&=\frac\ell2+\frac{i\pi}{4},\label{eq:atom1}\\
\Li_2(q)&=\frac{5\pi^2}{96}-\frac{\ell^2}{8}
 +i\left(G-\frac{\pi\ell}{8}\right),\label{eq:atom2}\\
\Rea\Li_3(q)&=\frac{35}{64}\zeta(3)
 -\frac{5\pi^2\ell}{192}+\frac{\ell^3}{48}.\label{eq:atom3}
\end{align}
We supply the short derivations to make the subsequent specialization
independent of a computer algebra system's branch convention.

\subsection{The required dilogarithm and trilogarithm evaluations}
Equation~\eqref{eq:atom1} is the principal logarithm.
The real part of \eqref{eq:atom2} follows from
\[
 \Li_2(q)+\Li_2(1-q)=\zeta(2)-\log q\log(1-q)
\]
and $1-q=\bar q$. The imaginary part follows from
\eqref{eq:Fr} with $r=1$ at $z=i$, since
$\Li_2(i)=-\pi^2/48+iG$; the latter equality follows directly by
separating even and odd terms in the absolutely convergent series.

For the trilogarithm use
\begin{align}
\Li_3(x)+\Li_3(1-x)+\Li_3\!\left(\frac{x}{x-1}\right)
={}&\zeta(3)+\zeta(2)\log(1-x)\notag\\
&-\frac12\log x\log^2(1-x)
 +\frac16\log^3(1-x).\label{eq:trilog}
\end{align}
On $0<x<1$ its derivative follows from the dilogarithm reflection
formula and
$\Li_2(x/(x-1))=-\Li_2(x)-\tfrac12\log^2(1-x)$.
The limit at $x=0$ fixes the integration constant. Continue into the
upper half-plane to $x=q$; no principal cut is crossed.
Here $q/(q-1)=-i$ and $\Rea\Li_3(-i)=-3\zeta(3)/32$.
Taking real parts in \eqref{eq:trilog} gives \eqref{eq:atom3}.

\subsection{The functional identities before specialization}
Theorem~\ref{thm:onezero} first gives the stronger complex formulas
\begin{align}
\Li_{2,1,1}(z)
={}&\Li_4(q)-L\Li_3(q)+\frac{L^2}{2}\Li_2(q)
-\frac{L^3}{6}\Li_1(q)\notag\\
&-\zeta(4)-\frac{L^4}{24},\label{eq:t211}\\
\Li_{1,2,1}(z)
={}&-3\Li_4(q)+2L\Li_3(q)-\frac{L^2}{2}\Li_2(q)
+L\zeta(3)\notag\\
&+3\zeta(4)-\frac{L^4}{24},\label{eq:t121}\\
\Li_{1,1,2}(z)
={}&3\Li_4(q)-L\Li_3(q)-2L\zeta(3)
-\frac{L^2}{2}\zeta(2)\notag\\
&-3\zeta(4)-\frac{L^4}{24}.\label{eq:t112}
\end{align}
Their full real parts are determined as well; only the imaginary parts
were printed in the source's candidate display.

\begin{theorem}[The three manuscript identities]\label{thm:gaussian}
All three of the following are exact:
\begin{align}
\Ima\Li_{2,1,1}(i)
={}&\lambda_4+\frac\ell2\lambda_3
-\frac{G(\pi^2-4\ell^2)}{32}
-\frac{\pi(8\ell^3+105\zeta(3))}{768},\label{eq:g211}\\
\Ima\Li_{1,2,1}(i)
={}&-3\lambda_4-\ell\lambda_3
+\frac{G(\pi^2-4\ell^2)}{32}
-\frac{3\pi^3\ell}{256}\notag\\
&+\frac{\pi(2\ell^3+67\zeta(3))}{128},\label{eq:g121}\\
\Ima\Li_{1,1,2}(i)
={}&3\lambda_4+\frac\ell2\lambda_3
+\frac{5\pi^3\ell}{192}-\frac{163\pi\zeta(3)}{256}.
\label{eq:g112}
\end{align}
\end{theorem}
\begin{proof}
Substitute $L=-\ell/2+i\pi/4$ and
\eqref{eq:atom1}--\eqref{eq:atom3} in
\eqref{eq:t211}--\eqref{eq:t112}, and collect imaginary parts.
No unproved relation among the $\lambda_n$ is used.
For example the only unknown real classical part that occurs before
simplification is $\Rea\Li_3(q)$, already evaluated in
\eqref{eq:atom3}; the coefficient of $\Li_4(q)$ is real.
The exact symbolic substitutions are reproduced by
\code{code/check_exact.py}.
\end{proof}

\subsection{The Eisenstein companion}
Equations~\eqref{eq:t211}--\eqref{eq:t112},
\eqref{eq:new5}, and~\eqref{eq:new6} apply unchanged at
\eqref{eq:eispoint}. Thus the same index families reduce to classical
polylogarithms at $q=1/2+i\sqrt3/6$, powers of
$L=-\tfrac12\log3+i\pi/6$, and ordinary zeta values.
This is already an explicit exact Eisenstein reduction; no independence
claim about those half-point classical values is needed.

For complete real and imaginary formulas define
$U_n=\Rea\Li_n(q)$ and $V_n=\Ima\Li_n(q)$ and use
\begin{equation}\label{eq:complexproduct}
\Ima\bigl(L^j\Li_n(q)\bigr)
=\Rea(L^j)V_n+\Ima(L^j)U_n.
\end{equation}
Every coefficient in \eqref{eq:complexproduct} is an explicit polynomial
in $\log3$ and $\pi$, with rational coefficients.
Keeping this small structured basket is preferable to asserting a
further reduction into a smaller root-of-unity basket without proof.

\section{A double-polylogarithm formula for an all-weight family}

The next formula addresses the first case beyond the classical ladder.
It reduces the height-one family with initial index $3$ to singles and
doubles at $q$, irrespective of its original depth.
Let
\begin{equation}\label{eq:Ck}
 C_1(q)=0,\qquad C_k(q)=\sum_{a=1}^{k-1}\Li_{a,k-a}(q)\quad(k\ge2).
\end{equation}
All MPLs on the right have the one-variable meaning \eqref{eq:mpl}.

\begin{theorem}[Two-zero height-one identity]\label{thm:twozero}
For every $w\ge3$ and $z\in D$,
\begin{align}
\Li_{3,\ones{w-3}}(z)
={}&\frac{L^w}{w!}
 +\sum_{k=1}^{w}\frac{(-1)^kL^{w-k}}{(w-k)!}
   \bigl((k-2)\Li_k(q)+C_k(q)\bigr)\notag\\
&+(-1)^{w-1}\Bigl((L+\Li_1(q))\zeta(w-1)
 +(w-2)\zeta(w)-\zeta(w-1,1)\Bigr).
\label{eq:twozero}
\end{align}
The right side has reduced depth at most two.
\end{theorem}

\begin{proof}
Write $D_q=q\,d/dq=d/dL$. Direct differentiation of the nested series,
followed by continuation, gives
\[
 D_qC_k=C_{k-1}+\frac{q}{1-q}\Li_{k-1}\quad(k\ge2).
\]
Put $V_k=(k-2)\Li_k+C_k$, with $V_1=-\Li_1$.
Then
\[
 D_qV_k=V_{k-1}+\frac{\Li_{k-1}}{1-q}\quad(k\ge2),
 \qquad D_qV_1=-\frac{q}{1-q}.
\]
Differentiating the finite sum in \eqref{eq:twozero} telescopes and yields
\[
 \frac{F_{w-2}(z)-(-1)^{w-2}\zeta(w-1)}{q-1}.
\]
The derivative of the explicit $(L+\Li_1(q))\zeta(w-1)$ term supplies
the missing constant term, because
$D_q(L+\Li_1(q))=1/(1-q)$.
The derivative of the entire right side is therefore
$F_{w-2}(z)/(q-1)$.
The left side has the same derivative, since
$d\Li_{3,\ones{w-3}}(z)/dz=F_{w-2}(z)/z$ and
$dz/z=dL/(q-1)$.

For the endpoint constant, recall the convergent depth-two sum identity
\begin{equation}\label{eq:eulersum}
 \sum_{a=2}^{w-1}\zeta(a,w-a)=\zeta(w),\qquad w\ge3.
\end{equation}
It follows, for example, by subtracting the shuffle and stuffle
expansions of the once-regularized product $\zeta(1)\zeta(w-1)$:
its common divergent term $\zeta(1,w-1)$ cancels, leaving precisely
\eqref{eq:eulersum}. Only one logarithmic divergence is involved, so
no higher regularization correction enters. Equivalently one may use
finite cutoffs and subtract the common harmonic-logarithmic part
before taking the limit.

The convergent colored stuffle identity, for $0<q<1$, gives
\[
 \Li_{1,w-1}(q)-\Li_1(q)\zeta(w-1)
 =-\Li_{w-1,1}(1,q)-\Li_w(q).
\]
The right side tends to $-\zeta(w-1,1)-\zeta(w)$ by dominated
convergence. Combining this with \eqref{eq:eulersum} yields
\[
 \lim_{q\to1^-}\bigl(C_w(q)-\Li_1(q)\zeta(w-1)\bigr)
 =-\zeta(w-1,1).
\]
All positive powers of $L$ multiplying at most logarithmically singular
terms vanish. The remaining constants in \eqref{eq:twozero} cancel,
so its right side tends to zero as $z\to0^-$. The left side does too.
Equality of derivatives and the endpoint limit proves the result;
continue to $D$.
\end{proof}

\begin{remark}[The sum identity and the proof dependency]
Equation~\eqref{eq:eulersum} is the classical depth-two Euler sum
formula, not a new conjecture assumed here. The same two-zero formula
can also be obtained directly from Theorem~\ref{thm:complement} by
expanding its two zero positions and normalizing trailing zeros.
The supplied code verifies the agreement of these two independently
specified formulas through weight twelve.
\end{remark}

At weight six, \eqref{eq:twozero} is a complete depth-two representation
of $\Li_{3,1,1,1}(i)$ and of $\Li_{3,1,1,1}(\rho)$.
In particular a depth-four index does not by itself identify a new
depth-four period in the enlarged comparison algebra.
The endpoint double is harmless as a named convergent constant;
standard Euler reductions can simplify it further. For example,
\[
 \zeta(4,1)=2\zeta(5)-\zeta(2)\zeta(3),\qquad
 \zeta(5,1)=\frac34\zeta(6)-\frac12\zeta(3)^2.
\]
These last simplifications are not required for the theorem.

\section{Exact formal shuffle ranks at every weight and depth}

\subsection{What is being counted}
Let $\mathcal H^1$ be the rational vector space with basis the empty
word and all binary words ending in $1$, equipped with the shuffle
product. Its positive-weight part is $\mathcal H^1_+$.
Define the indecomposable quotient
\[
 Q\mathcal H^1=\mathcal H^1_+/(\mathcal H^1_+)^2,
\]
where the square is taken with the shuffle product.
This is a formal quotient before evaluating any word at a complex point.
Weight is word length, and depth is the number of $1$ letters.

\begin{theorem}[Multigraded shuffle rank]\label{thm:rank}
For $1\le d\le w$, the weight-$w$, depth-$d$ part of
$Q\mathcal H^1$ has dimension
\begin{equation}\label{eq:witt}
 \ell_{w,d}=\frac1w\sum_{k\mid\gcd(w,d)}
 \mu(k)\binom{w/k}{d/k}.
\end{equation}
Consequently the space spanned by homogeneous shuffle products inside
the $\binom{w-1}{d-1}$ word coordinates has rank
\begin{equation}\label{eq:rank}
 R_{w,d}=\binom{w-1}{d-1}-\ell_{w,d}.
\end{equation}
A basis of the quotient is given by the binary Lyndon words of that
multidegree, with $0<1$.
\end{theorem}
\begin{proof}
The Chen--Fox--Lyndon factorization uniquely writes every word as a
concatenation of nonincreasing Lyndon words. The triangular shuffle
basis theorem \cite{Radford,MelanconReutenauer} says that the shuffle
product of those factors is
\[
 \left(\prod_\lambda m_\lambda!\right)w
    +\text{lexicographically smaller words},
\]
where $m_\lambda$ counts repeated factors.
A nontrivial binary Lyndon word begins in $0$ and ends in $1$.
In the factorization of a word ending in $1$, the one-letter factor
$0$ cannot occur: as the smallest factor it would occur last, making
the word end in $0$. Thus every factor is an allowed generator of
$\mathcal H^1$.

Triangularity gives a polynomial basis indexed by monomials in these
Lyndon generators. Passing to indecomposables leaves exactly the
single-generator monomials. This proves the asserted quotient basis,
without any evaluation or transcendence assumption.

For completeness, count its multidegree by primitive necklaces. There
are $\binom wd$ binary words of length $w$ containing $d$ ones.
Every word is uniquely a power of a primitive word. M\"obius inversion
therefore counts primitive words of this multidegree as
$\sum_{k\mid\gcd(w,d)}\mu(k)\binom{w/k}{d/k}$.
Each primitive rotation class contains exactly $w$ words and exactly
one Lyndon representative, proving \eqref{eq:witt}.
Finally, fixing the last letter to be $1$ gives
$\binom{w-1}{d-1}$ original coordinates, proving \eqref{eq:rank}.
\end{proof}

\subsection{The manuscript's ten triples}
For $d=3$ and $w\ge3$, \eqref{eq:witt} becomes
\begin{equation}\label{eq:triplecount}
 \ell_{w,3}=\frac{(w-1)(w-2)}6-\frac{\mathbf1_{3\mid w}}3.
\end{equation}
At $w=6$, the three Lyndon words are
\[
 000111,\quad001011,\quad001101,
\]
corresponding respectively to $(4,1,1)$, $(3,2,1)$, and $(3,1,2)$.
Thus $R_{6,3}=10-3=7$, exactly the reported formal count.

The accompanying file \code{weight6_triples_shuffle.json} supplies a
particularly transparent certificate: seven non-Lyndon factorization
products give a $7\times10$ rational matrix of rank seven. It is not
necessary to reconstruct the source's thirteen reported rows.
For example,
\begin{equation}\label{eq:sampleproduct}
 \Li_2(z)\Li_{3,1}(z)
 =9\Li_{4,1,1}(z)+4\Li_{3,2,1}(z)
 +\Li_{3,1,2}(z)+\Li_{2,3,1}(z).
\end{equation}
The coefficient of the non-Lyndon word $010011$ is one.

\begin{center}
\begin{tabular}{rrrr}
\toprule
Weight & Triple coordinates & Formal indecomposables & Product rank\\
\midrule
3&1&0&1\\
4&3&1&2\\
5&6&2&4\\
6&10&3&7\\
7&15&5&10\\
8&21&7&14\\
9&28&9&19\\
10&36&12&24\\
\bottomrule
\end{tabular}
\end{center}

\subsection{Why this is not a numerical depth theorem}
Evaluation at $z=i$ or $z=\rho$ is an algebra map out of the formal
shuffle algebra. It can have additional relations, and changing the
allowed argument family can introduce additional reductions.
Taking real or imaginary parts also changes how products are represented.
Therefore the three formal Lyndon coordinates at weight six need not be
three independent numbers. Theorems about parity give lower-depth
conjugation combinations, not independence of the other components
\cite{Panzer,Umezawa}.

The contrast with Section 4 is instructive. A one-zero multidegree has
one formal Lyndon generator, but that generator is nevertheless
classically reducible after the argument is changed to $q$. Formal
shuffle indecomposability is a property of a presentation, not a proof
that a numerical period has minimal depth.

\section{Corrections and the remaining \texorpdfstring{$S_4$}{S4} question}

\subsection{The mixed-point antisymmetry must swap colors}
The manuscript's discussion following the symmetric Gaussian doubles
says that the remaining content lies in
$\Li_{a,b}(P)-\Li_{b,a}(P)$ and that these differences vanish on the
diagonal. This is correct when both colors in $P$ are equal, but not
for a mixed point. The stuffle identity actually controls
\begin{equation}\label{eq:colorswap}
 \Li_{a,b}(x,y)+\Li_{b,a}(y,x)
 =\Li_a(x)\Li_b(y)-\Li_{a+b}(xy).
\end{equation}
The corresponding antisymmetric combination is
$\Li_{a,b}(x,y)-\Li_{b,a}(y,x)$.
Even when $a=b$, it need not vanish if $x\ne y$.

\begin{proposition}[An exact Gaussian diagonal counterexample]
\label{prop:mixedcorrection}
\begin{equation}\label{eq:counterexample}
 \Li_{1,1}(i,-1)-\Li_{1,1}(-1,i)
 =\frac{\pi^2}{24}-\frac{\log^22}{4}
 +i\left(G-\frac\pi2\log2\right)\ne0.
\end{equation}
\end{proposition}
\begin{proof}
Integrating the derivative in the first argument gives
\[
 \Li_{1,1}(i,-1)
 =\log2\log(1-i)-\Li_2\left(\frac{1-i}{2}\right)+\Li_2(1/2).
\]
This can be checked directly by differentiating the corresponding
formula with $i$ replaced by a variable $x$ along the segment from $0$
to $i$: the derivative is $-\log(1+x)/(1-x)$ and both sides vanish
at $x=0$. Principal branches are continuous along that segment.
Using \eqref{eq:atom2} and
$\Li_2(1/2)=\pi^2/12-\log^22/2$, we find
\[
 \Li_{1,1}(i,-1)
 =\frac{\pi^2}{32}+\frac{\log^22}{8}
 +i\left(G-\frac{3\pi}{8}\log2\right).
\]
On the other hand \eqref{eq:colorswap} gives the sum as
\[
 \frac{\log^22}{2}+\frac{\pi^2}{48}
 +i\left(G-\frac\pi4\log2\right).
\]
Twice the first value minus this sum proves \eqref{eq:counterexample}.
Its real part is positive: $\pi>3$ and $0<\log2<1$ already give a
lower bound greater than $1/8$.
\end{proof}

The correction is not that the manuscript's symmetric stuffle formulas
are wrong; they are correct. The error is the subsequent identification
of the involution on a mixed-color family. The integration notes supply
replacement wording and the counterexample.

\subsection{A precise finite obstruction for the attempted \texorpdfstring{$S_4$}{S4} proof}
Define $g_{a,b}=\Ima\Li_{a,b}(i,1)$ and
$S_p=\sum_{n\ge0}(-1)^nH_n/(2n+1)^p$.
An elementary parity filter gives, for $p>1$,
\begin{equation}\label{eq:Spfilter}
 S_p=g_{p,1}+\Ima\Li_{p,1}(i,-1).
\end{equation}
Indeed the sum of the inner colors $1+(-1)^m$ retains only even $m$,
and the imaginary part of $i^{2n+1}$ is $(-1)^n$.
Absolute convergence justifies the rearrangement.

Thus the printed $S_4$ formula is equivalent to
\begin{align}
3g_{4,1}+3g_{3,2}+9g_{2,3}
+7\Ima\Li_{4,1}(i,-1)
=\frac{\pi^5}{32}-\frac{27}{32}G\zeta(3)-14\beta(4)\log2.
\label{eq:S4target}
\end{align}
We attempted to derive it from all weight-five products of two
single polylogarithms with fourth-root colors, including the
once-regularized product differences and internal doubling distribution.
The finite test system has 134 rows and 23 imaginary-coordinate columns,
after conjugation and removal of divergent first-index-one coordinates.
Its rational rank is 19; adjoining the left side of
\eqref{eq:S4target} raises the rank to 20.

This is independently replayable by a particularly small dual witness.
Write $X_{a,b}^{x,y}=\Ima\Li_{a,b}(i^x,i^y)$, with conjugate columns
identified by sign. Give the following columns the listed rational
values and give every other column zero:
\begin{center}
\begin{tabular}{rrrrrrrr}
\toprule
$X_{1,4}^{1,0}$&$X_{2,3}^{0,1}$&$X_{2,3}^{1,0}$&$X_{2,3}^{1,3}$&
$X_{3,2}^{0,1}$&$X_{3,2}^{1,0}$&$X_{3,2}^{1,3}$&$X_{4,1}^{0,1}$\\
\midrule
$-2$&$1$&$3$&$-1$&$-3$&$-1$&$-1$&$2$\\
\bottomrule
\end{tabular}
\end{center}
Every homogeneous law row evaluates to zero, whereas the target left
side evaluates to $24$. The JSON file contains every column, every row,
and the witness; the script verifies $An=0$ and $v^Tn=24$ in exact
arithmetic.

This does \emph{not} disprove \eqref{eq:S4target}. It proves only that
this specified finite row space cannot supply the missing derivation.
Higher-depth relations, further argument transformations, or additional
cyclotomic functional equations can enlarge the row space. No decision
about the true period relation follows from the obstruction. We retain
$S_4$ as an unresolved identity in this contribution rather than
relabeling numerical agreement as proof.

\section{Exact computation, numerical diagnostics, and enclosures}

\subsection{The exact reduction engine}
The core file \code{code/polylog_words.py} implements
\eqref{eq:master}, the trailing-zero recursion, the one-zero closed
formula, the two-zero closed formula, and Lyndon factorization.
It uses Python integers and rational fractions for all coefficients.
The JSON normal form uses a key $(j,\boldsymbol a,\boldsymbol b)$ for
$L^j\Li_{\boldsymbol a}(q)\zeta(\boldsymbol b)$.
Empty indices mean one. Each output term is checked for weight and
reduced depth.

The finite verification run performed the following exact checks:
\begin{center}
\begin{tabular}{lr}
\toprule
Check & Count\\
\midrule
Trailing-zero normalization re-expanded in the shuffle algebra & 511\\
One-zero formula versus independent word transport & 66\\
Two-zero height-one formula versus word transport & 10\\
Lyndon triangular leading-word and multiplicity checks & 1023\\
Printed Gaussian identities after symbolic substitution & 3\\
High-depth catalogue entries, weight at most eight & 142\\
Terms in that catalogue & 5763\\
\bottomrule
\end{tabular}
\end{center}
In addition, the weight-six triple certificate is checked directly as a
$7\times10$ matrix of rational rank seven.
These finite checks test the implementations; the all-weight claims
come from the proofs, not extrapolation from the table.

\subsection{Independent quadrature checks}
The separate numerical implementation does not evaluate the original
side by invoking the word-transport formula. For one-zero words it
uses the single integral \eqref{eq:singleintegral}. For the two-zero
height-one family it uses
\begin{equation}\label{eq:twointegral}
 \Li_{3,\ones b}(z)=\frac1{(b+1)!}\int_0^1
 \frac{-\log t\,[-\log(1-zt)]^{b+1}}{t}\,dt.
\end{equation}
The transformed side uses convergent nested power series at $q$.
At 75 decimal digits, 153 cases at $z=-1/2$, $z=i$, and $z=\rho$
passed; the largest observed absolute residual was
$1.378\times10^{-73}$ (rounded upward for this display).
These mpmath quadratures are independent diagnostics, not rigorous
interval calculations.

\subsection{A simple rigorous series tail bound}
\begin{proposition}\label{prop:tail}
For $|q|=r<1$, positive indices $\boldsymbol s$ of depth $d$, and
$N\ge0$, let $M=N+1$. If
\[
 s_1\ge\frac{d-1}{1+\log M},
\]
then truncating the outer index of \eqref{eq:mpl} at $N$ gives error at
most
\begin{equation}\label{eq:tail}
 \frac{r^M(1+\log M)^{d-1}}
 {(1-r)(d-1)!M^{s_1}}.
\end{equation}
For $d=1$ this is $r^{N+1}/((1-r)(N+1)^{s_1})$ for every $N$.
\end{proposition}
\begin{proof}
The inner finite nested sum is at most
$H_{n-1}^{d-1}/(d-1)!$, because its indices are positive and the product
expansion of $H_{n-1}^{d-1}$ includes all strictly ordered distinct terms
$(d-1)!$ times. Also $H_{n-1}\le1+\log n$.
The logarithmic derivative of
$(1+\log x)^{d-1}/x^{s_1}$ is nonpositive for $x\ge M$ under the stated
condition. Bound this factor by its value at $M$ and sum the remaining
geometric series.
\end{proof}

\subsection{Rational interval certificates for the Gaussian triples}
The file \code{code/certify_gaussian.py} constructs rational endpoints,
not floating-point guesses. Classical polylogarithms at $(1+i)/2$ are
summed using exact rational real and imaginary parts for 850 terms.
Since $|(1+i)/2|<3/4$, each component's tail is bounded by
$4(3/4)^{851}/851^s$.
The constants $\pi$ and $\log2$ use, respectively, the alternating
Machin arctangent series and the positive $2\operatorname{artanh}(1/3)$
series. Catalan's constant is evaluated as
\[
 G=\frac1{16}\bigl(\zeta(2,1/4)-\zeta(2,3/4)\bigr).
\]
The Hurwitz zeta values and $\zeta(3)$ use Euler--Maclaurin summation
with rational Bernoulli numbers. For $s\ge2$, $x=N+a$, and $M$ retained
Bernoulli corrections, the standard periodic-Bernoulli remainder obeys
\begin{equation}\label{eq:EMbound}
 |R_M|\le
 \frac{|B_{2M}|}{(2M)!}(s)_{2M-1}x^{-s-2M+1}.
\end{equation}
This follows by bounding the periodic Bernoulli polynomial by
$|B_{2M}|$ in the integral remainder and integrating
$(x+t)^{-s-2M}$. The code uses $N=80$, $M=50$.

Interval arithmetic then propagates these enclosures through the proven
complex formulas and separately through the printed imaginary formulas.
All six target components have enclosure width less than $10^{-90}$;
the actual reported widths are smaller. The imaginary values begin
\begin{align*}
\Ima\Li_{2,1,1}(i)&=-0.0231000737339630776916090328202431130613391338545\ldots,\\
\Ima\Li_{1,2,1}(i)&=-0.0285334247842249902585459705838397251944563205761\ldots,\\
\Ima\Li_{1,1,2}(i)&=-0.0451614156780673253466795694168971191581439693134\ldots.
\end{align*}
The JSON output stores exact endpoint fractions as well as outward-rounded
decimal displays. These are \emph{proof-backed value enclosures}; overlap
of two intervals is not offered as a substitute for an analytic proof.

\section{Proposed further research}

\paragraph{The missing $S_4$ relation.}
Use the dual witness to select genuinely new equations rather than
repeating the same depth-two product system. The octahedral map
$t\mapsto(1-t)/(1+t)$ permutes the level-four punctures and is a concrete
candidate source. The first objective is an exact finite certificate
for \eqref{eq:S4target}, with all endpoint constants justified.

\paragraph{Weight-six triple candidates.}
Determine explicit reductions, or a carefully stated motivic obstruction,
for the imaginary components of the three Lyndon triples at $i$.
Specify whether half-point doubles and triples are allowed. A result in
one comparison algebra does not decide the question in another.

\paragraph{Minimal two-zero baskets.}
Formula~\eqref{eq:twozero} uses the symmetric index sums $C_k(q)$.
Investigate whether the entire two-zero sector can be organized into
fewer such combinations using shuffle and functional identities.
This is an exact elimination problem before any independence question.

\paragraph{Argument choice and conditioning.}
The map $q=(1-z)^{-1}$ is favorable at the two points studied here but
not uniformly over all $z$. Optimize among the six fractional-linear
permutations of $0,1,\infty$, tracking both the number of transformed
$1$ letters and the convergence radius. Include coefficient growth
and cancellation in the optimization criterion.

\paragraph{Arithmetic of the Eisenstein half-point.}
Compare the classical basket at $1/2+i\sqrt3/6$ with the fixed-level-three
cyclotomic basket. Prove any proposed bridge by functional equations
rather than relying on an enriched PSLQ search that reports a survivor.

\paragraph{Certified high-depth evaluation.}
Extend the rational or algebraic interval implementation from the
one-zero Gaussian family to the full catalogue. The difficult part is
certified MZV endpoint evaluation, not the exponentially convergent
$q$-series itself. An exact coefficient engine and an interval evaluator
should remain independent modules.

\paragraph{Formalization.}
A proof-assistant development can separate the finite word algebra from
analytic transport. The former needs finite words, shuffle, rational
normalization, and the Lyndon count. The latter needs branch domains,
regularized basepoints, and the finite-weight endpoint argument in
Lemma~\ref{lem:connection}. No such formalization is claimed in this
package.

\paragraph{Uniform relation catalogues.}
Replace isolated rank statements by public matrices, column conventions,
and dual or row-combination certificates. The all-weight count in
Theorem~\ref{thm:rank} is a model for the formal part; a corresponding
complete cyclotomic relation theorem would require additional laws and
must not be inferred from this binary-word calculation.

\section{Integration and reproducibility}

The package is additive: it does not overwrite the manuscript or commit
changes to the repository. The proposed integration location is a new
research-contribution directory under
\code{Analysis/Polylogarithms/docs/}, with the relevant theorem summaries
and correction copied into Chapter 4 after review.

\code{integration/EDITORIAL-NOTES.md} distinguishes a proven correction,
newly proved candidate displays, and unchanged open questions.
\code{integration/complementary-depth-insert.tex} gives a compact
manuscript-ready insert. The standalone article is the full proof source.

A local reproduction requires Python 3.10 or later, SymPy, mpmath, and a
standard LaTeX installation. Run
\begin{verbatim}
python code/check_exact.py
python code/check_numerical.py
python code/certify_gaussian.py
python code/audit_s4_rows.py
python code/replay_certificates.py
latexmk -pdf -interaction=nonstopmode -halt-on-error article.tex
\end{verbatim}
The numerical and interval runs are separate so that an exact-algebra
check never quietly substitutes for an analytic or numerical obligation.
The included manifest hashes the deliverable files. No original
repository file is silently edited, and no unsupported claim of Lean,
Wolfram, or independent peer verification is made.

\appendix
\section{The exact certificate formats}

Each reduction entry stores the original indices, weight, depth, and
an array of terms. A term with data
\begin{verbatim}
{"log_power": j, "q_indices": [a1,...],
 "zeta_indices": [b1,...],
 "numerator": p, "denominator": r}
\end{verbatim}
means $(p/r)L^j\Li_{a_1,\ldots}(q)\zeta(b_1,\ldots)$.
An empty index array means the factor $1$.
The coefficient is exact, and the denominator is positive.

\code{complementary_depth_catalogue.json} contains every word ending in
$1$ of weight at most eight whose depth is greater than half its weight:
142 entries and 5763 nonzero normal-form terms.
\code{one_zero.json} contains all 66 pairs $a,b$ of weights two through
twelve. Each formula applies at both special points and throughout $D$.
\code{two_zero_height_one.json} records ten identities of weights three
through twelve.

The word transport enumerates $2^K$ transformed words and $w+1$ splits
before normalization, so the raw number of connection terms is at most
$(w+1)2^K$. This is a count before trailing-zero expansion, \emph{not}
a polynomial-time claim for the full output. Normalization and the size
of rational coefficients can increase the computational cost. For a
fixed small zero count, the raw transport is particularly small, which
explains the practical success of the one- and two-zero sectors.

\section{The finite law system in the \texorpdfstring{$S_4$}{S4} audit}
For transparency, the rows are generated as follows. For every
$p+q=5$ with $p,q\ge1$ and $x,y\in\{1,i,-1,-i\}$, include the
imaginary parts of the stuffle and shuffle products when both single
factors converge. The shuffle is
\begin{align*}
\Li_p(x)\Li_q(y)
={}&\sum_{k=0}^{p-1}\binom{q-1+k}{k}
 \Li_{q+k,p-k}(y,x/y)\\
&+\sum_{k=0}^{q-1}\binom{p-1+k}{k}
 \Li_{p+k,q-k}(x,y/x).
\end{align*}
When one factor is $\Li_1(1)$, include only the once-regularized
shuffle-minus-stuffle difference, where the common divergent coordinate
cancels. Add the internal doubling rows
\[
 \sum_{\epsilon,\eta\in\{\pm1\}}
 \Li_{a,b}(\epsilon x,\eta y)
 =2^{2-a-b}\Li_{a,b}(x^2,y^2),
\]
using $x,y\in\{1,i\}$ and omitting rows with an uncancelled divergent
coordinate. Some rows are zero or duplicate; they are retained so that
the generator prescription remains simple and reproducible.

The dual witness is a vector in the kernel of the matrix of these
homogeneous left sides. Its pairing with the target is nonzero.
This certificate does not require or imply that the right sides of all
possible cyclotomic identities have been exhausted.

\clearpage
\begin{thebibliography}{99}
\raggedright
\bibitem{ProveIt}
V. Reshetnikov, \emph{ProveIt, Analysis/Polylogarithms/docs/manuscript},
Chapters 4 and 10, repository snapshot
\code{29d9344771b41df39ff9b7890c76b1f152572dde}, inspected 9 October 2026.
\url{https://github.com/VladimirReshetnikov/ProveIt/tree/29d9344771b41df39ff9b7890c76b1f152572dde/Analysis/Polylogarithms/docs/manuscript}.

\bibitem{RemiddiVermaseren}
E. Remiddi and J. A. M. Vermaseren,
\emph{Harmonic polylogarithms}, International Journal of Modern Physics A
\textbf{15} (2000), 725--754.
\url{https://arxiv.org/abs/hep-ph/9905237}.

\bibitem{Kolbig}
K. S. K\"olbig, \emph{Nielsen's generalized polylogarithms},
SIAM Journal on Mathematical Analysis \textbf{17} (1986), 1232--1258.
DOI: 10.1137/0517086.
\url{https://repository.cern/records/09vyq-j8461}.

\bibitem{Radford}
D. E. Radford, \emph{A natural ring basis for the shuffle algebra and an
application to group schemes}, Journal of Algebra \textbf{58} (1979),
432--454. DOI: 10.1016/0021-8693(79)90171-6.

\bibitem{MelanconReutenauer}
G. Melan\c{c}on and C. Reutenauer,
\emph{Lyndon words, free algebras and shuffles}, Canadian Journal of
Mathematics \textbf{41} (1989), 577--591.
\url{https://reutenauer.math.uqam.ca/wp-content/uploads/2024/04/LyndonWordsFreeAlgebras-Shuffles1989.pdf}.

\bibitem{Panzer}
E. Panzer, \emph{The parity theorem for multiple polylogarithms},
Journal of Number Theory \textbf{172} (2017), 93--113.
\url{https://arxiv.org/abs/1512.04482}.

\bibitem{Umezawa}
R. Umezawa, \emph{An explicit parity theorem for multiple polylogarithms},
preprint, 2025. \url{https://arxiv.org/abs/2508.02040}.
\end{thebibliography}
\end{document}
